\documentclass[aspectratio=1610,12pt,notheorems]{beamer}

\usepackage[utf8x]{inputenc} \usepackage[russian]{babel}
\usepackage{amsmath,amssymb,amsthm,mathtools}
\usepackage{graphicx,caption,subcaption}
\usepackage{hyperref,natbib}
\usepackage{tikz,xcolor,colortbl,makecell}
\usepackage{algorithm,algpseudocode}
\usepackage{qrcode}

\usetikzlibrary{arrows,backgrounds,patterns,%
	matrix,shapes,fit,calc,shadows,plotmarks,snakes}

\theoremstyle{plain}
\newtheorem{theorem}{Теорема}
\newtheorem{lemma}[theorem]{Лемма}

\theoremstyle{definition}
\newtheorem{definition}{Определение}
\newtheorem{problem}{Задача}

\usetheme[height=0.97cm]{Rochester}
\usecolortheme{dolphin}

\definecolor{hard}{RGB}{145,55,55}
\definecolor{mnsgold}{RGB}{240,235,220}

\setbeamercolor{headline}{bg=hard,fg=mnsgold}
\setbeamercolor*{frametitle}{parent=headline}

\setbeamercolor{structure}{fg=hard}
\setbeamercolor{subsection in head/foot}{bg=white,fg=hard}
\setbeamercolor{section in head/foot}{bg=hard,fg=mnsgold}
\setbeamercolor{block title}{bg=hard,fg=mnsgold}

\setbeamertemplate{navigation symbols}{}

\parskip=2.6mm
\def\ll{\left(} \def\rr{\right)}
\def\lag{\left\langle} \def\rag{\right\rangle}

\definecolor{failpos}{RGB}{230,30,20}
\definecolor{initpos}{RGB}{30,20,220}
\definecolor{turna}{RGB}{100,240,110}
\definecolor{turnb}{RGB}{250,140,110}

\newcommand{\vseper}{\vphantom{$\int_{0_0}^{0^0}$}}

\newcommand{\singlepayoff}[2]{\tikz{
	\draw (0,0) node{\vseper #1}; \draw (0.8,0.8) node{\vseper #2};
}}

\newcommand{\rowdescription}[1]{\tikz{
	\draw (0,0) node{\vseper #1}; \draw (0,0.8) node{\vseper};
}}

\newcommand{\coldescription}[1]{\tikz{
	\draw (0,0) node{\vseper}; \draw (0.8,0) node{\vseper #1};
}}

\newcommand{\myref}[2]{\href{#1}{\texttt{\underline{#2}}}}

\def\mitem{\medskip\item}
\def\ps{\\ [0.65cm]} \linespread{1.16}
\def\fram#1#2{\begin{frame}\frametitle{#1}#2\end{frame}}
\def\usl#1#2{\begin{block}{#1} #2 \end{block} \medskip\pause}
\def\uslnp#1#2{\begin{block}{#1} #2 \end{block} \medskip}
\def\mov#1#2{\begin{scope}[xshift = #1 cm] #2 \end{scope};}

\def\divsby{\mathrel{\rlap{.}\rlap{\raisebox{0.55ex}{.}}\raisebox{1.1ex}{.}}}
\definecolor{starvert}{RGB}{110,175,230}
\newcommand{\vfi}{\varphi}
\newcommand{\litem}{\vspace{0.5cm}\item}

\newcommand{\nkstar}[2]{
	\foreach \i in {1,...,#1} {\draw[thick] (90 + 360 * \i / #1 : 3.5)
		-- (90 + 360 * \i / #1 + 360 * #2 / #1 : 3.5);}
	\foreach \i in {1,...,#1} {\fill[starvert] (90 + 360 * \i / #1 : 3.5) circle[radius=0.16cm];}
}

\title[Solutions of MNS]
    {\bfseries Решения избранных задач}

\author[\ ]
	{Е.И. Тодоров,\quad «Математика НОН-СТОП»\\ \vspace{0.3cm}
		{\small Фонд «Время Науки»}}

\institute[\ ]{\ }

\date{\today}

%%%%%%%%%%%%%%
%%%%%%%%%%%%%%

\begin{document}

\frame{\titlepage}


\begin{frame} \frametitle{Эта презентация онлайн}
\begin{center}
\vspace{-10mm}
Зачем фотографировать презентацию, когда её можно скачать?\vspace{3mm}

		\qrcode[height=4cm]{http://bit.ly/3_razbor_ei}\vspace{3mm}
		
		\url{http://bit.ly/3_razbor_ei}
\end{center}
\end{frame}

\begin{frame}
	\ \\ [-0.5cm]
	\noindent В 2020–2021 гг. олимпиада «Математика НОН-СТОП»  и\\
	 Санкт-Петербургский Tурнир юных математиков проводятся \\
	с использованием гранта Президента Российской Федерации \\
	на развитие гражданского общества, предоставленного \\
	{\bfseries Фондом президентских грантов.}
	
	\vspace{0.45cm}
	\begin{center}
		\includegraphics[scale=0.6]{fpg/fpg}
	\end{center}
\end{frame}

%%%%%% 4 класс, кубики
\begin{frame} \frametitle{Кубики}
\begin{figure}[h] \centering
    \tikz[scale=0.47]{
    {
    \draw[black!50!white,thick]{ (-4,2)--(4,2)--(4,4)--(-4,4)--(-4,2)
    (-2,2)--(-2,6)--(0,6)--(0,0)--(2,0)--(2,4)
    }
    }
    %6
    {
    {\fill[black] (-1.4,5.5) circle (0.15cm)}
    {\fill[black] (-0.6,5.5) circle (0.15cm)}
    {\fill[black] (-1.4,5) circle (0.15cm)}
    {\fill[black] (-0.6,5) circle (0.15cm)}
    {\fill[black] (-1.4,4.5) circle (0.15cm)}
    {\fill[black] (-0.6,4.5) circle (0.15cm)}
    }
    %5
    {
    {\fill[black] (-3,3) circle (0.15cm)}
    {\fill[black] (-3.4,3.4) circle (0.15cm)}
    {\fill[black] (-3.4,2.6) circle (0.15cm)}
    {\fill[black] (-2.6,3.4) circle (0.15cm)}
    {\fill[black] (-2.6,2.6) circle (0.15cm)}
    }
    %4
    {
    {\fill[black] (-1.4,3.4) circle (0.15cm)}
    {\fill[black] (-1.4,2.6) circle (0.15cm)}
    {\fill[black] (-0.6,3.4) circle (0.15cm)}
    {\fill[black] (-0.6,2.6) circle (0.15cm)}
    }
    %2
    {
    {\fill[black] (1.4,3.4) circle (0.15cm)}
    {\fill[black] (0.6,2.6) circle (0.15cm)}
    }
    %3
    {
    {\fill[black] (3.4,3.4) circle (0.15cm)}
    {\fill[black] (3,3) circle (0.15cm)}
    {\fill[black] (2.6,2.6) circle (0.15cm)}
    }
    %1
    {
    {\fill[black] (1,1) circle (0.15cm)} } }
\end{figure}


	\usl{2020-4-6B}{Сколько видимых точек может быть на башне из $6$ кубиков?}

Сумма чисел на противоположных гранях равна 7. \medskip \pause

Ответ — $84 + k$.
\end{frame}


%%%%%% 4 класс, разрезания

\begin{frame} \frametitle{Кирпичей требуют наши сердца}

\usl{2020-4-4C}{
	Нарисуйте на клетчатой бумаге такую фигуру, которую
	можно разделить по клеткам на 2, на 3, на 4, на 5, на 6
	одинаковых по форме и размеру связных фигур~—
	причём они не будут прямоугольниками.
}

\begin{center} \tikz[scale=0.28]
{
   \foreach \x / \y in {
	1/1, 1/2, 1/3, 1/4, 1/5, 1/6, 1/7, 1/8, 1/9, 1/10,  %делим на 2
	2/2, 2/3, 2/4, 2/5, 2/6, 2/7, 2/8, 2/9, 2/10, 2/11,
	3/1, 3/2, 3/3, 3/4, 3/5, 3/6, 3/7, 3/8, 3/9, 3/10,
	9/1, 9/2, 9/3, 9/4, 9/5, 9/6, 9/7, 9/8, 9/9, 9/10,  %делим на 3
	10/2, 10/3, 10/4, 10/5, 10/6, 10/7, 10/8, 10/9, 10/10, 10/11,
	13/1, 13/2, 13/3, 13/4, 13/5, 13/6, 13/7, 13/8, 13/9, 13/10,
	14/2, 14/3, 14/4, 14/5, 14/6, 14/7, 14/8, 14/9, 14/10, 14/11,
	17/1, 17/2, 17/3, 17/4, 17/5,    %делим на 4
	18/2, 18/3, 18/4, 18/5, 18/6,
	19/1, 19/2, 19/3, 19/4, 19/5, 
	20/7, 20/8, 20/9, 20/10, 20/11,
	21/6, 21/7, 21/8, 21/9, 21/10,
	22/7, 22/8, 22/9, 22/10, 22/11,
	25/3, 25/4, 25/7, 25/8,     %делим на 5
	26/4, 26/5, 26/8, 26/9, 
	27/3, 27/4, 27/7, 27/8, 
	28/4, 28/5, 28/8, 28/9, 
	29/3, 29/4, 29/7, 29/8, 
	30/4, 30/5, 30/8, 30/9, 
	33/1, 33/2, 33/3, 33/4, 33/5,     %делим на 6
	34/2, 34/3, 34/4, 34/5, 34/6, 
	35/6, 35/7, 35/8, 35/9, 35/10,
	36/7, 36/8, 36/9, 36/10, 36/11,
	37/1, 37/2, 37/3, 37/4, 37/5, 
	38/2, 38/3, 38/4, 38/5, 38/6, 38/7
   } 
   \draw[fill opacity=0.6] (\x, \y) rectangle (\x cm + 1 cm, \y cm + 1 cm);
   \foreach \x / \y in {
	4/2, 4/3, 4/4, 4/5, 4/6, 4/7, 4/8, 4/9, 4/10, 4/11,    %делим на 2
	5/1, 5/2, 5/3, 5/4, 5/5, 5/6, 5/7, 5/8, 5/9, 5/10,
	6/2, 6/3, 6/4, 6/5, 6/6, 6/7, 6/8, 6/9, 6/10, 6/11,
	11/1, 11/2, 11/3, 11/4, 11/5, 11/6, 11/7, 11/8, 11/9, 11/10,    %делим на 3
	12/2, 12/3, 12/4, 12/5, 12/6, 12/7, 12/8, 12/9, 12/10, 12/11,
	17/6, 17/7, 17/8, 17/9, 17/10,    %делим на 4
	18/7, 18/8, 18/9, 18/10, 18/11,
	19/6, 19/7, 19/8, 19/9, 19/10,
	20/2, 20/3, 20/4, 20/5, 20/6, 
	21/1, 21/2, 21/3, 21/4, 21/5, 
	22/2, 22/3, 22/4, 22/5, 22/6,
	25/1, 25/2, 25/5, 25/6, 25/9, 25/10,     %делим на 5
	26/2, 26/3, 26/6, 26/7, 26/10, 26/11,
	27/1, 27/2, 27/5, 27/6, 27/9, 27/10,
	28/2, 28/3, 28/6, 28/7, 28/10, 28/11,
	29/1, 29/2, 29/5, 29/6, 29/9, 29/10,
	30/2, 30/3, 30/6, 30/7, 30/10, 30/11,
	33/6, 33/7, 33/8, 33/9, 33/10,     %делим на 6
	34/7, 34/8, 34/9, 34/10, 34/11,
	35/1, 35/2, 35/3, 35/4, 35/5,
	36/2, 36/3, 36/4, 36/5, 36/6,
	37/6, 37/7, 37/8, 37/9, 37/10,
	38/8, 38/9, 38/10, 38/11
   } 
   \draw[fill=gray,fill opacity=0.6] (\x, \y) rectangle (\x cm + 1 cm, \y cm + 1 cm);
   \node at (2.5,6.5) {1};     %нумеруем 2
   \node at (5.5,6.5) {2};
   \node at (9.5,6.5) {1};     %нумеруем 3
   \node at (11.5,6.5) {2};
   \node at (13.5,6.5) {3};
   \node at (18.5,9.5) {1};     %нумеруем 4
   \node at (21.5,8.5) {2};
   \node at (18.5,4.5) {3};
   \node at (21.5,3.5) {4};
   \node at (27.5,9.5) {1};     %нумеруем 5
   \node at (27.5,7.5) {2};
   \node at (27.5,5.5) {3};
   \node at (27.5,3.5) {4};
   \node at (27.5,1.5) {5};
   \node at (33.5,9.5) {1};     %нумеруем 6
   \node at (35.5,9.5) {2};
   \node at (37.5,9.5) {3};
   \node at (33.5,4.5) {4};
   \node at (35.5,4.5) {5};
   \node at (37.5,4.5) {6};
} 
\end{center}
\end{frame}

\begin{frame} \frametitle{Гном вернулся --- дома нет}

\uslnp{2021-4-3A}{
Старый дом находится в лесной деревне Домовиково, старый гном отошёл по делам в деревню Гномовиково. Между деревнями 1 километр. Одновременно, соскучившись друг по другу, гном выходит из Гномовиково и дом выходит из Домовиково. Они ходят между деревнями туда и обратно, пока не окажутся одновременно в одной деревне — они не останавливаются, даже если встретят друг друга по пути.\\
Дом ходит со скоростью $3\text{ км/ч}$, гном ходит со скоростью $4\text{ км/ч}$. Сколько раз они встретятся по дороге между Домовиково и Гномовиково, пока не окажутся одновременно в одной деревне?
}
\end{frame}

\begin{frame} \frametitle{Гном вернулся --- дома нет}
   \begin{center} \tikz[scale=0.9]{
	\draw (0,0) rectangle (6,6);
	\draw (0,0) node[below]{Г} node[left]{0\,ч} (6,0) node[below]{Д} (0,6) node[left]{1\,ч};
	\draw[thick] (0,0) -- (6,1.5) -- (0,3) -- (6,4.5) -- (0,6)
		-- (6,4) -- (0,2) -- (6,0);
   } \end{center}
\end{frame}

%%%%%% 5 класс, разрезания
\begin{frame} \frametitle{Разрезания}
	\usl{2020-5-1C}{
		Можно ли нарисовать на клетчатом листе бумаги такую фигуру, которую можно разрезать по линиям сетки на две {\itshape одинаковые} фигуры двумя способами~— причём фигуры в первом и во втором способе были бы одни и те же, но линии разреза выглядели бы по-разному?
	}

\begin{center} \tikz[scale=0.75]
{
   \foreach \x / \y in {
	1/3, 2/3, 3/3, 3/4, 6/1, 6/2, 6/3, 7/3
   } 
	\draw[thick,fill opacity=0.6] (\x, \y) rectangle (\x cm + 1 cm, \y cm + 1 cm);
   {
   \foreach \x / \y in {
	1/1, 1/2, 2/2, 3/2, 7/2, 8/2, 8/3, 8/4
   } 
	\draw[thick,fill=gray,fill opacity=0.6] (\x, \y) rectangle (\x cm + 1 cm, \y cm + 1 cm);
	\node at (1.5,2.5) {$A$};
	\node at (3.5,3.5) {$B$};
	\node at (6.5,3.5) {$B$};
	\node at (8.5,2.5) {$A$};
}
} 
\end{center}
\end{frame}


%%%%%% 5 класс, семнашка, 4 пары

\begin{frame}
\frametitle{Аксиомы выборов}

\uslnp{2019-7-3A}{На предприятии работают 50 человек, и они выбирают себе начальника. Есть две кандидатуры, Ваня и Даня. Про каждого работника известно заранее, кому он отдаёт предпочтение: 20 человек за Даню, 30 человек за Ваню. \smallskip \\
Голосование проходит по двухтуровой системе: люди делятся на 5 групп по 10 человек, в каждой группе выбирается кандидат, наиболее популярный среди членов этой группы, и затем из 5 ответов выбирается имя, названное большее число раз. \smallskip \\
Разделите работников на группы так, чтобы в большинстве групп выбрали Даню и он победил на выборах, несмотря на изначально меньшее число голосующих за него.}
\end{frame}

\begin{frame}
\frametitle{Аксиомы выборов}

\begin{center} \tikz{
\begin{scope}[rotate=90,scale=1.2]
	\filldraw[fill=gray,draw=gray,opacity=0.32] (0,0) rectangle (1,5);
	\foreach \x in {0,2,5} {\draw[thick, color=gray]
		(0.5 * \x cm, 0) -- (0.5 * \x cm, 5);}
	\foreach \x in {0,10} {\draw[thick, color=gray]
		(0, 0.5 * \x cm) -- (2.5, 0.5 * \x cm);}
	\foreach \x in {1,3,4} {\draw[color=gray, opacity=0.38]
		(0.5 * \x cm, 0) -- (0.5 * \x cm, 5);}
	\foreach \x in {1,...,9} {\draw[color=gray, opacity=0.38]
		(0, 0.5 * \x cm) -- (2.5, 0.5 * \x cm);}

	\draw (0.5,-0.25) node[right]{\large За Даню};
	\draw (1.75,-0.25) node[right]{\large За Ваню};

	\pause
	\begin{scope} [xscale=-1, xshift=-2.5cm]
	\draw[very thick] (0.5,0) -- (0.5,5);
	\draw[very thick] (2.5,1.5) -- ++ (-1,0) --
		++ (0,0.5) -- ++ (-0.5,0) -- ++(0,-2);
	\draw[very thick] (2.5,3.5) -- ++ (-1,0) --
		++ (0,-0.5) -- ++ (-0.5,0) -- ++(0,2);
	\draw[very thick] (2.5,3.5) --
		++ (-1,0) -- ++ (0,-0.5) --
		++ (-0.5,0) -- ++ (0,-1) -- ++ (0.5,0) --
		++ (0,-0.5) -- ++(1,0);
	\end{scope}
\end{scope} } \end{center} \end{frame}

\begin{frame} \frametitle{Полиминное поле}
	\uslnp{2021-5-5A}{
	Покажите, как покрыть доску $8\times8$, используя 21 прямоугольник $1\times3$ (или $3\times1$) и один квадратик $1\times1$.
	}
\begin{center} \tikz[scale=0.55]{
   \foreach \x / \y in {
	1/1, 1/2, 1/6, 1/7, 1/8,
	3/4, 3/5,
	6/1, 6/2, 6/6, 6/7, 6/8
   } {
	\draw[thick,fill opacity=0.6] (\x, \y) rectangle (\x cm + 3 cm, \y cm + 1 cm);
   }
   \foreach \x / \y in {
	1/3,
	2/3,
	4/1,4/6,
	5/1,5/6,
	6/3,
	7/3,
	8/3
   } {
	\draw[thick,fill opacity=0.6] (\x, \y) rectangle (\x cm + 1 cm, \y cm + 3 cm);
   }
} \end{center}
\end{frame}

\begin{frame} \frametitle{Полиминное поле}
	\usl{2021-5-5B}{
	Докажите, что в предыдущем пункте квадратик $1\times1$ нельзя ставить в угол поля.
	}

\begin{center} \tikz[scale=0.55]{
   \foreach \x / \y in {
	1/1,
	1/4, 2/3, 3/2, 4/1,
	1/7, 2/6, 3/5, 4/4, 5/3, 6/2, 7/1,
	3/8, 4/7, 5/6, 6/5, 7/4, 8/3,
	6/8, 7/7, 8/6
   } {
		\draw[thick,fill=gray, fill opacity=0.6] (\x, \y) rectangle (\x cm + 1 cm, \y cm + 1 cm);
   }
   \foreach \x / \y in {
	1/2, 2/1,
	1/5, 2/4, 3/3, 4/2, 5/1,
	1/8, 2/7, 3/6, 4/5, 5/4, 6/3, 7/2, 8/1,
	4/8, 5/7, 6/6, 7/5, 8/4,
	7/8, 8/7
   } {
		\draw[thick,fill=black, fill opacity=0.6] (\x, \y) rectangle (\x cm + 1 cm, \y cm + 1 cm);
   }
   \foreach \x / \y in {
	1/3, 1/6,
	2/8, 
	3/1,
	5/8, 
	6/1, 
	8/2, 8/5, 8/8
   } {
		\draw[thick, fill opacity=0.6] (\x, \y) rectangle (\x cm + 1 cm, \y cm + 1 cm);
   }
   \node at (8.5,8.5) {\Large $\star$};
} \end{center}
\pause
Белых --- 20, серых --- 21, чёрных --- 22.
\end{frame}



%%%%%% 6 класс, научки
\begin{frame} \frametitle{Опубликовать за 60 секунд}
	\usl{2020-6-2B}{
	Научные руководители придумывают темы работ.
     \begin{itemize}
	\item Один из них придумывает 1 новую тему;
	\item После этого кто-то из них придумывает 2 новых темы;
	\item После этого кто-то из них придумывает 3 новых темы.
     \end{itemize}
	Пусть изначально первый придумал на $n$ тем больше, чем второй. Докажите, что руководители всегда смогут сравнять количество придуманных ими тем.
	}

Каждые два хода разность количеств тем будем сокращать на 1.
\end{frame}

\begin{frame}\frametitle{Роботы играют в карты}
\uslnp{2021-6-6C}{
Робот $f$ берёт натуральное число $n$, переставляет его цифры задом наперёд его и избавляется от ведущих нулей. Например, если положить в него число 2010, то $f$ вернёт число 102, будем записывать это как $f(2010) = 102$. Аналогично, $f(123) = 321$ и $f(14000) = 41$. \\
Робот $g$ берёт карточку с натуральным числом $n$, а возвращает карточку с числом $\frac{n}{f\left(f(n)\right)}$. Например, $g(100500) = \frac{100500}{f\left(f(100500)\right)} = \frac{100500}{f(5001)} = \frac{100500}{1005} = 100$. Олег складывает в робота $g$ карточки от 1 до некоторого числа $n$. Как количество различных карточек, возвращаемых роботом, будет зависеть от $n$? Что это будут за карточки?
}
\end{frame}

\begin{frame}\frametitle{Роботы играют в карты}
\begin{enumerate}
	\item Если у $n$ на конце нет нулей, то $f\left(f(n)\right) = n$, то есть $g(n) = 1$. \pause
	\item Если же $n$ делится на $10$, то $n = m\cdot10^k$, где $m$ --- число без нулей на конце. Из предыдущего пункта мы знаем, что $f\left(f(n)\right) = m$. Тогда $$g(n) = \frac{ m\cdot10^k}{m} = 10^k.$$ \pause
	\item Значит, различных значений будет ровно столько, сколько различных степеней 10 содержится среди чисел от 1 до $n$: $\lfloor \log_{10}n \rfloor$.
\end{enumerate}
\end{frame}

\begin{frame}\frametitle{}
\usl{2021-7-9A}{
Определим операцию $\star$ для положительных чисел $a$ и $b$ следующим образом: $a\star b=\frac{ab}{a+b}$.\\
Докажите, что $(a\star b)\star c=a\star (b\star c)$.
}
\usl{2021-7-9A}{
Посчитайте $1\star\left(2\star\left(4\star\left(\ldots\star\left(256\star(512\star1024)\right)\ldots\right)\right)\right)$.
}
Перерасставим скобки: $$(\ldots((1\star2)\star4)\star\ldots)\star1024 = \left(\ldots\left(\frac{2}{3}\star4\right)\star\ldots\right)\star1024 = \left(\ldots\frac{4}{7}\star\ldots\right)\star1024 = \frac{1024}{2047}.$$
\end{frame}

%%%%%% Профиль
\begin{frame} \frametitle{Профильные задания}
{\it Система високосных лет} для числа $t$ — это последовательность натуральных чисел $(a_0, a_2, a_3, \ldots, a_n)$ такая, что $a_{i+1}$ делится на $a_i$, а также
	$$\frac{1}{a_0} - \frac{1}{a_1} + \frac{1}{a_2} - \frac{1}{a_3} + \ldots
	     + (-1)^{n} \cdot \frac{1}{a_n}\ =\ t.$$ \smallskip

Какой могла бы быть система високосных лет, если бы длина года составляла 365.21875, 365.17, 365.33 дней? Для любого ли рационального числа существует система високосных лет?
\end{frame}

\begin{frame} \frametitle{Профильные задания}

\begin{align*}
	& 365.21875 = 365 + \frac{1}{4} - \frac{1}{32}
		\vphantom{\int\limits_{0_0}^{0^0}} \\
	& 365.17 = 365 + \frac{1}{5} - \frac{1}{25} + \frac{1}{100}
		\vphantom{\int\limits_{0_0}^{0^0}} \\
	& 365.33 = 365 + \frac{1}{3} - \frac{1}{300}
		\vphantom{\int\limits_{0_0}^{0^0}} \\
\end{align*}

\end{frame}

\begin{frame} \frametitle{Летний Турнир юных математиков}
\begin{itemize}
	\item нестандартные задачи;
	\item готовят к формату Турнира;
	\item 3 серии по 3--4 задачи;
	\item на решение даётся по 1,5--2 дня;
	\item Турнир проводится очно со 2 по 8 августа;
	\item решения проверяют авторы задач.
\end{itemize}
\end{frame}

\begin{frame} \frametitle{Летний Турнир юных математиков}
\begin{center}
\vspace{-10mm}
Собирайте команду и вступайте в бой!\vspace{3mm}

		\qrcode[height=4cm]{https://forms.gle/5fAsDNECis7aFZKT8}\vspace{3mm}
		
		\url{spbtym.ru} $\rightarrow$ Зарегистрироваться
\end{center}
\end{frame}

\begin{frame} \begin{center}
\ 

	{\Large Спасибо за внимание!} \vspace{0.4cm}

	{\large \url{mathnonstop.ru}}

	{\large \url{spbtym.ru}}

	{\large \url{mathnonstop@timeforscience.ru}}

	{\large \url{vk.com/timeforscience}}
\end{center} \end{frame}

\end{document}
